\documentclass{article}
\usepackage[french]{babel}
\usepackage[T1]{fontenc}
\usepackage[utf8]{inputenc}
\usepackage{graphicx}
\usepackage{mathtools}
\usepackage{listingsutf8}
\usepackage{amsfonts}
\usepackage{amssymb}
\usepackage{xcolor}
\lstset{%
    language=Python,
    inputencoding=utf8,
    basicstyle=\ttfamily\small,
    keywordstyle=\bfseries,
    showstringspaces=false,
    backgroundcolor=\color{gray!10},
    literate=
        {é}{{\'e}}{1}%
        {è}{{\`e}}{1}%
        {à}{{\`a}}{1}%
        {â}{{\^a}}{1}%
        {ç}{{\c{c}}}{1}%
        {œ}{{\oe}}{1}%
        {ù}{{\`u}}{1}%
        {É}{{\'E}}{1}%
        {È}{{\`E}}{1}%
        {À}{{\`A}}{1}%
        {Ç}{{\c{C}}}{1}%
        {Œ}{{\OE}}{1}%
        {Ê}{{\^E}}{1}%
        {ê}{{\^e}}{1}%
        {î}{{\^i}}{1}%
        {ï}{{\"i}}{1}%
        {ô}{{\^o}}{1}%
        {û}{{\^u}}{1}%
}
\usepackage{fullpage}
\usepackage[margin=2.1cm]{geometry}
\usepackage{amsthm}

\newtheoremstyle{exostyle}
{10pt}{10pt}{}{}
{\scshape\bfseries\large}
{\hfill\vspace{5pt}\newline}{0pt}
{\hfill\thmname{#1}\thmnumber{ #2} -- \thmnote{ #3}}

\newtheoremstyle{partiestyle}
{1em}{1em}{}{}{\bfseries}
{\vspace{.5em}\newline}{0em}
{\thmnumber{#2} \thmnote{ #3}}

\newtheoremstyle{questionstyle}
{.5em}{.5em}{}{}{\bfseries}
{}{0em}{Question \thmnumber{#2 }}

\theoremstyle{exostyle}
\newtheorem{exo}{Exercice}

\theoremstyle{partiestyle}
\newtheorem{partie}{}[exo]

\theoremstyle{questionstyle}
\newtheorem{questionpartie}{Question}[partie]

\newenvironment{reponse}{%
    \par\smallskip\noindent\textbf{Réponse.}\enspace\ignorespaces
}{%
    \par\smallskip
}

\usepackage{booktabs}
\usepackage{array}

\title{Corrigé -- Examen terminal: Programmation}
\author{L1 MPCI}
\date{28 mai 2026}

\begin{document}

\maketitle

% ============================================================
\section*{Barème (total : 20 points)}

\begin{center}
\begin{tabular}{llr}
\toprule
\textbf{Question} & \textbf{Contenu} & \textbf{Points} \\
\midrule
\multicolumn{2}{l}{\textit{Exercice 1 — Programmation orientée objet}} & \textbf{8,5} \\
\midrule
1.1.1 & Attributs d'instance (liste + types) & 0,5 \\
1.1.2 & Encapsulation (\texttt{\_}) + méthode \texttt{est\_disponible} & 0,5 \\
1.1.3 & Méthode \texttt{emprunter} & 0,5 \\
1.1.4 & Méthode \texttt{\_\_str\_\_} & 0,5 \\
1.2.1 & \texttt{ajouter} + \texttt{livres\_disponibles} & 0,5 \\
1.2.2 & \texttt{chercher} + \texttt{emprunter} (Bibliothèque) & 1,5 \\
1.3.1 & Définition de l'héritage & 0,5 \\
1.3.2 & Classe \texttt{LivreNumérique} + \texttt{super()} & 0,5 \\
1.3.3 & Surcharge \texttt{emprunter} & 0,5 \\
1.3.4 & Surcharge \texttt{\_\_str\_\_} & 0,5 \\
1.3.5 & Polymorphisme (justification) & 1 \\
1.4.1 & Boîte UML de \texttt{Livre} & 0,5 \\
1.4.2 & \texttt{LivreNumérique} + flèche d'héritage & 0,5 \\
1.4.3 & \texttt{Bibliothèque} + relation avec \texttt{Livre} + multiplicité & 0,5 \\
\midrule
\multicolumn{2}{l}{\textit{Exercice 2 — Séparation des responsabilités}} & \textbf{9,5} \\
\midrule
2.1.1 & Classification M/V/C des 6 éléments (0,5 pt chacun) & 3 \\
2.1.2 & Deux inconvénients du code monolithique & 1 \\
2.1.3 & Parties à réécrire pour une interface web & 1 \\
2.2.1 & Différence héritage / composition + exemples & 1 \\
2.2.2 & Constructeur de \texttt{Application} & 0,5 \\
2.2.3 & Méthode \texttt{emprunter} dans \texttt{Application} & 0,5 \\
2.3.1 & Diagramme UML \texttt{Application}/\texttt{Bibliothèque} + nature de la relation & 1,5 \\
2.3.2 & Créer vs recevoir les dépendances (composition vs association) & 1 \\
\midrule
\multicolumn{2}{l}{\textit{Exercice 3 — Implémentation du design pattern MVC}} & \textbf{8} \\
\midrule
3.1.1 & \texttt{VueTerminal} : 3 méthodes (0,5 pt chacune) & 1,5 \\
3.2.1 & \texttt{Contrôleur} : \texttt{\_\_init\_\_}, \texttt{afficher\_disponibles}, \texttt{emprunter} & 2 \\
3.2.2 & \texttt{Contrôleur.lancer} & 1 \\
3.3.1 & Code de lancement de l'application & 0,5 \\
3.3.2 & Méthodes requises pour \texttt{VueGraphique} + principe de substitution & 1 \\
3.4.1 & Diagramme UML complet MVC & 1 \\
3.4.2 & Ajout \texttt{VueGraphique} + principe de substitution & 1 \\
\bottomrule
\textbf{Total} & & \textbf{26} \\
\bottomrule
\end{tabular}
\end{center}

\clearpage

% ============================================================
\begin{exo}[Programmation orientée objet]

    \begin{partie}[La classe \texttt{Livre}]

        \begin{questionpartie}
            Qu'appelle-t-on un \textit{attribut d'instance}~? Listez les attributs d'instance de la classe \texttt{Livre} et donnez leur type.
        \end{questionpartie}
        \begin{reponse}
            Un attribut d'instance est une variable propre à chaque objet, créée dans \texttt{\_\_init\_\_} via \texttt{self}. Chaque instance possède sa propre copie indépendante de ces attributs. La classe \texttt{Livre} en possède trois~:
            \begin{itemize}
                \item \texttt{titre} : \texttt{str}
                \item \texttt{auteur} : \texttt{str}
                \item \texttt{\_disponible} : \texttt{bool}
            \end{itemize}
        \end{reponse}

        \begin{questionpartie}
            Pourquoi l'attribut \texttt{\_disponible} est-il préfixé d'un tiret bas~? Ajoutez une méthode \texttt{est\_disponible(self) -> bool}.
        \end{questionpartie}
        \begin{reponse}
            Le tiret bas indique par convention que l'attribut est \textit{protégé}~: il ne doit pas être lu ni modifié directement depuis l'extérieur de la classe. On passe par des méthodes dédiées pour respecter l'encapsulation.
\begin{lstlisting}
def est_disponible(self) -> bool:
    return self._disponible
\end{lstlisting}
        \end{reponse}

        \begin{questionpartie}
            Complétez \texttt{emprunter(self) -> bool}.
        \end{questionpartie}
        \begin{reponse}
\begin{lstlisting}
def emprunter(self) -> bool:
    if self._disponible:
        self._disponible = False
        return True
    return False
\end{lstlisting}
        \end{reponse}

        \begin{questionpartie}
            Écrivez \texttt{\_\_str\_\_(self) -> str}.
        \end{questionpartie}
        \begin{reponse}
\begin{lstlisting}
def __str__(self) -> str:
    etat = "disponible" if self._disponible else "emprunté"
    return f"{self.titre} ({self.auteur}) [{etat}]"
\end{lstlisting}
        \end{reponse}

    \end{partie}

    \begin{partie}[La classe \texttt{Bibliothèque}]

        \begin{questionpartie}
            Ajoutez \texttt{ajouter} et \texttt{livres\_disponibles}.
        \end{questionpartie}
        \begin{reponse}
\begin{lstlisting}
def ajouter(self, livre: Livre):
    self._livres.append(livre)

def livres_disponibles(self) -> list:
    return [l for l in self._livres if l.est_disponible()]
\end{lstlisting}
        \end{reponse}

        \begin{questionpartie}
            Ajoutez \texttt{chercher} puis \texttt{emprunter}.
        \end{questionpartie}
        \begin{reponse}
\begin{lstlisting}
def chercher(self, titre: str):
    for livre in self._livres:
        if livre.titre == titre:
            return livre
    return None

def emprunter(self, titre: str) -> bool:
    livre = self.chercher(titre)
    if livre is None:
        return False
    return livre.emprunter()
\end{lstlisting}
        \end{reponse}

    \end{partie}

    \begin{partie}[Héritage : le livre numérique]

        \begin{questionpartie}
            Qu'appelle-t-on l'héritage~?
        \end{questionpartie}
        \begin{reponse}
            L'héritage permet à une classe \textit{fille} de reprendre automatiquement les attributs et méthodes d'une classe \textit{parente}. La classe fille est une spécialisation de la parente (relation \og est-un \fg{})~: tout objet de la classe fille \textit{est aussi} un objet de la classe parente. La classe fille peut ajouter de nouveaux membres ou \textit{surcharger} (redéfinir) ceux hérités.
        \end{reponse}

        \begin{questionpartie}
            Écrivez la classe \texttt{LivreNumérique} héritant de \texttt{Livre}.
        \end{questionpartie}
        \begin{reponse}
\begin{lstlisting}
class LivreNumérique(Livre):
    def __init__(self, titre: str, auteur: str):
        super().__init__(titre, auteur)
\end{lstlisting}
            L'appel à \texttt{super().\_\_init\_\_()} délègue l'initialisation des attributs (\texttt{titre}, \texttt{auteur}, \texttt{\_disponible}) à la classe parente.
        \end{reponse}

        \begin{questionpartie}
            Surchargez \texttt{emprunter} pour qu'elle retourne toujours \texttt{True}.
        \end{questionpartie}
        \begin{reponse}
\begin{lstlisting}
def emprunter(self) -> bool:
    return True
\end{lstlisting}
            On ne modifie pas \texttt{\_disponible}~: un livre numérique est toujours disponible, son état ne change jamais.
        \end{reponse}

        \begin{questionpartie}
            Surchargez \texttt{\_\_str\_\_}.
        \end{questionpartie}
        \begin{reponse}
\begin{lstlisting}
def __str__(self) -> str:
    return f"{self.titre} ({self.auteur}) [numérique]"
\end{lstlisting}
        \end{reponse}

        \begin{questionpartie}
            Peut-on ajouter un \texttt{LivreNumérique} à une \texttt{Bibliothèque}~? Quel principe illustre ce comportement~?
        \end{questionpartie}
        \begin{reponse}
            Oui. \texttt{Bibliothèque.ajouter} attend un \texttt{Livre}~; or \texttt{LivreNumérique} hérite de \texttt{Livre}, donc tout \texttt{LivreNumérique} \textit{est} un \texttt{Livre} et peut le remplacer partout. Lorsque \texttt{Bibliothèque.emprunter} appellera \texttt{livre.emprunter()}, Python exécutera la version surchargée de \texttt{LivreNumérique}, qui retourne toujours \texttt{True}. Ce mécanisme s'appelle le \textbf{polymorphisme}~: un même appel de méthode produit un comportement différent selon le type réel de l'objet.
        \end{reponse}

    \end{partie}

    \begin{partie}[Modélisation UML]

        \begin{questionpartie}
            Dessinez la boîte UML de la classe \texttt{Livre}.
        \end{questionpartie}
        \begin{reponse}
\begin{verbatim}
+-----------------------------------+
|              Livre                |
+-----------------------------------+
| - titre : str                     |
| - auteur : str                    |
| - _disponible : bool              |
+-----------------------------------+
| + est_disponible() : bool         |
| + emprunter() : bool              |
| + __str__() : str                 |
| + __init__(titre, auteur)         |
+-----------------------------------+
\end{verbatim}
        \end{reponse}

        \begin{questionpartie}
            Ajoutez \texttt{LivreNumérique} au diagramme.
        \end{questionpartie}
        \begin{reponse}
\begin{verbatim}
+-------------------------------+
|         LivreNumérique        |
+-------------------------------+
|                               |
+-------------------------------+
| + emprunter() : bool          |
| + __str__() : str             |
+-------------------------------+
            △  (héritage)
            |
+-----------------------------------+
|              Livre                |
+-----------------------------------+
| ...                               |
+-----------------------------------+
\end{verbatim}
            La boîte de \texttt{LivreNumérique} ne répète pas les membres hérités~; on n'y inscrit que les méthodes surchargées. L'attribut \texttt{\_disponible} n'apparaît pas car aucun attribut nouveau n'est ajouté.
        \end{reponse}

        \begin{questionpartie}
            Ajoutez \texttt{Bibliothèque} au diagramme.
        \end{questionpartie}
        \begin{reponse}
\begin{verbatim}
+-------------------------------+        1..*    +-------+
|         Bibliothèque          | ------------> | Livre |
+-------------------------------+               +-------+
| - _livres : list              |
+-------------------------------+
| + ajouter(livre)              |
| + chercher(titre) : Livre     |
| + emprunter(titre) : bool     |
| + livres_disponibles() : list |
+-------------------------------+
\end{verbatim}
            La relation est une \textbf{composition}~: une \texttt{Bibliothèque} contient (et gère la durée de vie de) ses livres. La multiplicité \texttt{1..*} indique qu'une bibliothèque peut contenir un ou plusieurs livres (on peut aussi écrire \texttt{0..*} si une bibliothèque vide est admise).
        \end{reponse}

    \end{partie}
\end{exo}

% ============================================================
\begin{exo}[Séparation des responsabilités]

    \begin{partie}[Analyse]

        \begin{questionpartie}
            Classifiez les éléments \texttt{(a)} à \texttt{(f)} selon M, V ou C.
        \end{questionpartie}
        \begin{reponse}
            \begin{itemize}
                \item \texttt{(a)} \texttt{choix = input(...)} \textrightarrow{} \textbf{Vue}~: c'est une saisie utilisateur.
                \item \texttt{(b)} \texttt{dispo = [l for l in livres if l["dispo"]]} \textrightarrow{} \textbf{Modèle}~: filtrage des données métier.
                \item \texttt{(c)} \texttt{print(f"- \{l['titre']\}...")} \textrightarrow{} \textbf{Vue}~: affichage d'information à l'écran.
                \item \texttt{(d)} \texttt{titre = input("Titre : ")} \textrightarrow{} \textbf{Vue}~: saisie utilisateur.
                \item \texttt{(e)} \texttt{livre = next(...)} \textrightarrow{} \textbf{Modèle}~: recherche dans les données.
                \item \texttt{(f)} \texttt{livre["dispo"] = False} \textrightarrow{} \textbf{Modèle}~: modification de l'état d'un livre.
            \end{itemize}
            Note~: les blocs \texttt{if/elif} qui décident quelle action effectuer selon \texttt{choix} relèvent du \textbf{Contrôleur}.
        \end{reponse}

        \begin{questionpartie}
            Donnez deux inconvénients du code monolithique.
        \end{questionpartie}
        \begin{reponse}
            \begin{enumerate}
                \item \textbf{Non réutilisabilité}~: impossible de changer d'interface (terminal, web, graphique) sans réécrire la totalité de la fonction, y compris la logique métier.
                \item \textbf{Difficulté de maintenance}~: toute correction ou évolution oblige à lire l'intégralité du code mélangé, avec un risque fort d'introduire des régressions dans des parties sans lien avec la modification.
            \end{enumerate}
        \end{reponse}

        \begin{questionpartie}
            Que faudrait-il réécrire pour passer à une interface web~?
        \end{questionpartie}
        \begin{reponse}
            Toutes les parties \textbf{Vue}~: les \texttt{print(...)} (remplacés par du HTML/JSON) et les \texttt{input(...)} (remplacés par des requêtes HTTP). Dans une architecture MVC correcte, seule la Vue change~; le Modèle (règles métier, données) et le Contrôleur (logique de décision) restent intacts.
        \end{reponse}

    \end{partie}

    \begin{partie}[Refactoring par composition]

        \begin{questionpartie}
            Différence entre héritage et composition.
        \end{questionpartie}
        \begin{reponse}
            \begin{itemize}
                \item \textbf{Héritage} (relation \og est-un \fg{})~: \texttt{LivreNumérique} est un \texttt{Livre}, il hérite de ses attributs et méthodes. On l'utilise quand la classe fille est une spécialisation de la parente.
                \item \textbf{Composition} (relation \og a-un \fg{})~: \texttt{Bibliothèque} contient des \texttt{Livre}, elle ne les étend pas. On l'utilise quand un objet est composé d'autres objets sans en être une variante.
            \end{itemize}
        \end{reponse}

        \begin{questionpartie}
            Constructeur de \texttt{Application} par composition.
        \end{questionpartie}
        \begin{reponse}
\begin{lstlisting}
class Application:
    def __init__(self):
        self._bibliotheque = Bibliothèque()
\end{lstlisting}
            \texttt{Application} ne hérite pas de \texttt{Bibliothèque}~; elle en \textit{possède} une instance. C'est la composition.
        \end{reponse}

        \begin{questionpartie}
            Méthode \texttt{emprunter(self)} de \texttt{Application}.
        \end{questionpartie}
        \begin{reponse}
\begin{lstlisting}
def emprunter(self):
    titre = input("Titre : ")
    succes = self._bibliotheque.emprunter(titre)
    if succes:
        print("Emprunté.")
    else:
        print("Indisponible ou introuvable.")
\end{lstlisting}
            La logique métier (l'emprunt) est maintenant déléguée à \texttt{Bibliothèque}~; la méthode ne manipule plus directement les données.
        \end{reponse}

    \end{partie}

    \begin{partie}[Modélisation UML]

        \begin{questionpartie}
            Dessinez le diagramme de classes UML de \texttt{Application} ainsi que son lien avec \texttt{Bibliothèque}.
        \end{questionpartie}
        \begin{reponse}
\begin{verbatim}
+-------------------------------+          +-------------------------------+
|         Application           | <>-----> |         Bibliothèque          |
+-------------------------------+   1    1 +-------------------------------+
| - _bibliotheque : Bibliothèque|          | - _livres : list              |
+-------------------------------+          +-------------------------------+
| + emprunter()                 |          | + ajouter(livre)              |
+-------------------------------+          | + chercher(titre)             |
                                           | + emprunter(titre) : bool     |
                                           | + livres_disponibles() : list |
                                           +-------------------------------+
\end{verbatim}
            La relation est une \textbf{composition} (losange plein, \texttt{<>}) car \texttt{Application} \textit{crée} la \texttt{Bibliothèque} dans son constructeur et en est propriétaire : si l'application disparaît, la bibliothèque disparaît avec elle. La multiplicité \texttt{1\,:\,1} signifie qu'une \texttt{Application} possède exactement une \texttt{Bibliothèque}. Une simple \textbf{association} (flèche sans losange) s'utiliserait si les deux objets avaient des cycles de vie indépendants.
        \end{reponse}

        \begin{questionpartie}
            Créer vs recevoir les dépendances en UML.
        \end{questionpartie}
        \begin{reponse}
            \begin{itemize}
                \item \textbf{Contrôleur crée ses dépendances}~: relation de \textbf{composition} (losange plein) de \texttt{Contrôleur} vers \texttt{Bibliothèque} et vers \texttt{VueTerminal}. Le Contrôleur possède ces objets.
                \item \textbf{Contrôleur reçoit ses dépendances}~: relation d'\textbf{association} simple (flèche sans losange). Les objets existent indépendamment du contrôleur.
            \end{itemize}
            Le second cas est préférable pour la testabilité~: on peut passer au contrôleur une fausse vue ou un faux modèle (\textit{mock}) lors des tests unitaires, sans modifier le code du contrôleur.
        \end{reponse}

    \end{partie}
\end{exo}

% ============================================================
\begin{exo}[Implémentation du \textit{design pattern} MVC]

    \begin{partie}[La Vue]

        \begin{questionpartie}
            Implémentez en Python la classe \texttt{VueTerminal}.
        \end{questionpartie}
        \begin{reponse}
\begin{lstlisting}
class VueTerminal:
    def afficher_livres(self, livres: list):
        if not livres:
            print("Aucun livre disponible.")
        else:
            for livre in livres:
                print(str(livre))

    def afficher_message(self, message: str):
        print(message)

    def demander_titre(self) -> str:
        return input("Titre : ")
\end{lstlisting}
            La Vue ne contient aucune logique métier~: elle se contente d'afficher et de transmettre les saisies.
        \end{reponse}

    \end{partie}

    \begin{partie}[Le Contrôleur]

        \begin{questionpartie}
            Implémentez \texttt{\_\_init\_\_}, \texttt{afficher\_disponibles} et \texttt{emprunter}.
        \end{questionpartie}
        \begin{reponse}
\begin{lstlisting}
class Contrôleur:
    def __init__(self, modele: Bibliothèque, vue: VueTerminal):
        self._modele = modele
        self._vue = vue

    def afficher_disponibles(self):
        livres = self._modele.livres_disponibles()
        self._vue.afficher_livres(livres)

    def emprunter(self):
        titre = self._vue.demander_titre()
        succes = self._modele.emprunter(titre)
        if succes:
            self._vue.afficher_message("Emprunté avec succès.")
        else:
            self._vue.afficher_message("Indisponible ou introuvable.")
\end{lstlisting}
            Le Contrôleur ne fait jamais d'affichage direct (pas de \texttt{print}) ni d'accès direct aux données~: il délègue toujours à la Vue ou au Modèle.
        \end{reponse}

        \begin{questionpartie}
            Implémentez \texttt{lancer}.
        \end{questionpartie}
        \begin{reponse}
\begin{lstlisting}
    def lancer(self):
        while True:
            print("1. Livres disponibles  2. Emprunter  0. Quitter")
            choix = input("Votre choix : ")
            if choix == "1":
                self.afficher_disponibles()
            elif choix == "2":
                self.emprunter()
            elif choix == "0":
                break
\end{lstlisting}
            On pourrait déplacer le \texttt{print} et l'\texttt{input} du menu dans la Vue (méthodes \texttt{afficher\_menu} et \texttt{demander\_choix}) pour une séparation encore plus stricte.
        \end{reponse}

    \end{partie}

    \begin{partie}[Extensibilité]

        \begin{questionpartie}
            Écrivez le code de lancement de l'application.
        \end{questionpartie}
        \begin{reponse}
\begin{lstlisting}
modele = Bibliothèque()
modele.ajouter(Livre("Le Petit Prince", "Antoine de Saint-Exupéry"))
modele.ajouter(Livre("1984", "George Orwell"))

vue = VueTerminal()
controleur = Contrôleur(modele, vue)
controleur.lancer()
\end{lstlisting}
        \end{reponse}

        \begin{questionpartie}
            Quelles méthodes \texttt{VueGraphique} doit-elle posséder~?
        \end{questionpartie}
        \begin{reponse}
            \texttt{VueGraphique} doit exposer exactement les mêmes méthodes que celles appelées par \texttt{Contrôleur} sur la vue~:
            \begin{itemize}
                \item \texttt{afficher\_livres(self, livres: list)}
                \item \texttt{afficher\_message(self, message: str)}
                \item \texttt{demander\_titre(self) -> str}
            \end{itemize}
            C'est le principe de \textbf{substitution}~: \texttt{VueGraphique} peut remplacer \texttt{VueTerminal} partout sans que \texttt{Contrôleur} ne soit modifié, car il n'interagit avec la vue qu'au travers de ces trois méthodes. Pour passer à une interface graphique, seule la ligne \texttt{vue = VueTerminal()} changerait en \texttt{vue = VueGraphique()}~; le reste du code de lancement est inchangé.
        \end{reponse}

    \end{partie}

    \begin{partie}[Modélisation UML]

        \begin{questionpartie}
            Diagramme de classes complet de l'architecture MVC.
        \end{questionpartie}
        \begin{reponse}
\begin{verbatim}
+-------------------------------+
|         Bibliothèque          |
+-------------------------------+    (Modèle)
| - _livres : list              |
+-------------------------------+
| + ajouter(livre)              |
| + chercher(titre)             |
| + emprunter(titre) : bool     |
| + livres_disponibles() : list |
+-------------------------------+
             ^  (association)
             | 1
+-------------------------------+
|          Contrôleur           |
+-------------------------------+
| - _modele : Bibliothèque      |    1
| - _vue : VueTerminal          | ------->  +-------------------------------+
+-------------------------------+           |         VueTerminal           |
| + afficher_disponibles()      |           +-------------------------------+
| + emprunter()                 |           |                               |
| + lancer()                    |           +-------------------------------+
+-------------------------------+           | + afficher_livres(livres)     |
                                            | + afficher_message(msg)       |
                                            | + demander_titre() : str      |
                                            +-------------------------------+
\end{verbatim}
            \texttt{Contrôleur} est associé à \texttt{Bibliothèque} et à \texttt{VueTerminal} (associations simples, sans losange~: il reçoit ses dépendances en paramètre). La multiplicité est \texttt{1} des deux côtés~: un contrôleur utilise exactement un modèle et une vue.
        \end{reponse}

        \begin{questionpartie}
            Ajout de \texttt{VueGraphique} et principe de substitution.
        \end{questionpartie}
        \begin{reponse}
\begin{verbatim}
+-------------------------------+
|         VueTerminal           |  (Vue concrète 1)
+-------------------------------+
| + afficher_livres(livres)     |
| + afficher_message(msg)       |
| + demander_titre() : str      |
+-------------------------------+

+-------------------------------+
|         VueGraphique          |  (Vue concrète 2)
+-------------------------------+
| + afficher_livres(livres)     |
| + afficher_message(msg)       |
| + demander_titre() : str      |
+-------------------------------+

Contrôleur -------> VueTerminal   (ou VueGraphique)
\end{verbatim}
            \texttt{VueGraphique} n'hérite pas de \texttt{VueTerminal}~: ce sont deux classes indépendantes qui exposent la même \textit{interface} (mêmes signatures de méthodes). En UML strict on modéliserait une interface \texttt{IVue} dont les deux héritent, mais en Python le typage canard suffit.

            Le diagramme illustre le \textbf{principe de substitution}~: \texttt{Contrôleur} ne connaît que les méthodes \texttt{afficher\_livres}, \texttt{afficher\_message} et \texttt{demander\_titre}. N'importe quel objet proposant ces trois méthodes peut remplacer la vue sans modifier \texttt{Contrôleur}.
        \end{reponse}

    \end{partie}
\end{exo}

\end{document}
